\chapter{Exact orbital rotations at finite interaction}
\label{chap:orbitalrotations}

Chapter~\ref{chap:generalhopping} showed that rotating only the one-body
matrix is incomplete at finite $U$.  We now implement the missing two-body
term and turn basis invariance into an executable test.

\section{One convention for every index}

Let $R$ be unitary and define the new orbital operators by
\begin{equation}
c_{i\sigma}=\sum_a R_{ia}d_{a\sigma}.
\label{eq:rotationconvention}
\end{equation}
The columns of $R$ are therefore the new orbitals expressed in the old basis.
The one-body matrix becomes
\begin{equation}
h'_{ab}=\sum_{ij}R^*_{ia}h_{ij}R_{jb}
       =\left(R^\dagger hR\right)_{ab}.
\end{equation}

For a general up--down interaction, use the explicit operator convention
\begin{equation}
H_V=\sum_{ijkl}V_{ijkl}
c_{i\uparrow}^\dagger c_{j\uparrow}
c_{k\downarrow}^\dagger c_{l\downarrow}.
\end{equation}
Substitution of Eq.~\eqref{eq:rotationconvention} gives
\begin{equation}
V'_{abcd}=\sum_{ijkl}
R^*_{ia}R_{jb}R^*_{kc}R_{ld}V_{ijkl}.
\label{eq:generaltensorrotation}
\end{equation}
For the site-local Hubbard tensor $V_{iiii}=U$, this reduces to
Eq.~\eqref{eq:rotatedinteraction}.  Hermiticity is expressed directly in the
stored index order:
\begin{equation}
V_{abcd}=V^*_{badc}.
\end{equation}

The public transformation functions have separate, explicit roles:
\begin{itemize}
\item \code{rotate\_one\_body} transforms $h$;
\item \code{rotate\_interaction\_tensor} transforms $V$;
\item \code{rotate\_integrals} transforms a consistent $(h,V)$ pair;
\item \code{rotate\_hubbard\_integrals} starts from local $U$.
\end{itemize}
Inputs are copied, checked for finite values, checked for Hermiticity or
unitarity to tolerance $10^{-12}$, and returned read-only.

\section{Fermionic action and the strongest small-system test}

The tensor Hamiltonian applies
$c_{a\uparrow}^\dagger c_{b\uparrow}
 c_{c\downarrow}^\dagger c_{d\downarrow}$ from right to left using the same
global spin-orbital convention as every earlier operator.  It supports both
CSR construction and matrix-free products.  The local tensor has only $L$
nonzero entries, while a generic rotation makes all $L^4$ entries nonzero.

A unitary basis change cannot alter the exact many-body eigenvalues.  For a
small sector we can therefore compare the \emph{complete} sorted spectra,
including degeneracies, rather than checking only the ground energy.  This
single test probes the rotation conjugations, all four tensor indices,
fermionic signs, and matrix assembly.

\lstinputlisting[caption={Exact finite-$U$ orbital rotation and full-spectrum check.},label={lst:orbitalrotation}]{code/step08_orbital_rotation.py}

The listing also exposes the cost change: a local $U$ tensor becomes dense.
The reference implementation intentionally retains every nonzero coefficient
and estimates attempted action terms before accepting the calculation.  Near
half filling, a dense rotated tensor is quick at $L=4$, useful but slower at
$L=6$, and rejected by the default work guard at $L=8$.  Structured or
truncated tensors will require separate error controls in future work.

\section{What an observable means after rotation}

The spectrum is basis invariant, but an array index in the rotated basis labels
an orbital, not a physical site.  Applying the existing local-charge routine
directly to a rotated state measures orbital occupation.  To recover physical
site density or spin correlations, transform the observable tensor as well, or
transform the many-body state back.  This distinction is essential when
comparing wavelets, natural orbitals, or learned rotations.

\paragraph{Checkpoint.}
Use a deterministic complex unitary.  Verify $R^\dagger R=I$, Hermiticity of
the transformed arrays, complete-spectrum agreement below $10^{-10}$, CSR
versus matrix-free agreement, and small eigenpair residuals.  As a negative
control, rotate $h$ but retain local $U$ and confirm that the finite-$U$
spectrum changes.
